\documentclass[11pt]{article}
\usepackage[a4paper,margin=1in]{geometry}
\usepackage{fontspec}
\usepackage[boldfont=false]{xeCJK}
\setCJKmainfont{FandolSong-Regular.otf}
\usepackage{amsmath}
\usepackage{amssymb}
\usepackage{array}
\usepackage{booktabs}
\usepackage{graphicx}
\usepackage{adjustbox}
\usepackage{enumitem}
\setlist[itemize]{topsep=2pt,partopsep=0pt,parsep=0pt,itemsep=2pt,leftmargin=1.4em}
\setlist[enumerate]{topsep=2pt,partopsep=0pt,parsep=0pt,itemsep=2pt,leftmargin=1.6em}
\usepackage{parskip}
\usepackage{fvextra}
\fvset{breaklines=true,breakanywhere=true,fontsize=\small}
\setlength{\parindent}{0pt}

\begin{document}

\begin{center}
  {\LARGE\bfseries The Periodic Table}\\[0.35em]
  {\large Igcse Chemistry}
\end{center}
\vspace{0.6em}

\section*{Arrangement of the elements}

The \textbf{Periodic Table} 周期表 arranges all the \textbf{elements} 元素 in order of increasing \textbf{proton number} 质子数 (the \textbf{atomic number} 原子序数). The horizontal rows are called \textbf{periods} 周期 and the vertical columns are called \textbf{groups} 族.

A few key patterns:

\begin{itemize}
\item Across a period, the elements change from \textbf{metals} 金属 on the left to \textbf{non-metals} 非金属 on the right.

\item The group number tells you the charge of the \textbf{ions} 离子 that the elements form. Group I forms $+1$ ions, Group II forms $+2$ ions, and Group VII forms $-1$ ions.

\item Elements in the same group have similar chemical properties. This is because they have the same number of outer-shell \textbf{electrons} 电子 (the same outer \textbf{electronic configuration} 电子排布).

\end{itemize}

Because of these patterns, you can use an element's position to predict its properties.

\section*{Group I — the alkali metals}

Group I elements are the \textbf{alkali metals} 碱金属: \textbf{lithium} 锂, \textbf{sodium} 钠 and \textbf{potassium} 钾. They are \textbf{soft} 柔软 metals (you can cut them with a knife).

Going down the group, there are clear \textbf{trends} 趋势:

\begin{center}
\adjustbox{max width=\linewidth}{%
\begin{tabular}{ll}
\toprule
\textbf{Property} & \textbf{Trend going down the group} \\
\midrule
\textbf{melting point} 熔点 & decreases \\
\textbf{density} 密度 & increases \\
\textbf{reactivity} 活泼性 & increases \\
\bottomrule
\end{tabular}}
\end{center}

You can use these trends to predict the properties of other Group I elements. For example, rubidium (below potassium) would be even more reactive and have an even lower melting point.

\section*{Group VII — the halogens}

Group VII elements are the \textbf{halogens} 卤素: \textbf{chlorine} 氯气, \textbf{bromine} 溴 and \textbf{iodine} 碘. They are \textbf{diatomic} 双原子 non-metals (each molecule is made of two atoms, such as $\text{Cl}_2$).

Their appearance at room temperature and pressure:

\begin{center}
\adjustbox{max width=\linewidth}{%
\begin{tabular}{ll}
\toprule
\textbf{Halogen} & \textbf{Appearance} \\
\midrule
chlorine & a pale yellow-green gas \\
bromine & a red-brown liquid \\
iodine & a grey-black solid \\
\bottomrule
\end{tabular}}
\end{center}

Going down the group, the density increases but the reactivity \textbf{decreases} (the opposite trend to Group I).

\subsection*{Displacement reactions}

A more reactive halogen pushes out (displaces) a less reactive \textbf{halide} 卤化物 ion from its solution. This is a \textbf{displacement reaction} 置换反应.

For example, chlorine is more reactive than bromine, so chlorine displaces bromine from potassium bromide:

\[
\text{Cl}_2 + 2\text{KBr} \rightarrow 2\text{KCl} + \text{Br}_2
\]

\section*{Transition elements}

The \textbf{transition elements} 过渡元素 are the block of metals in the middle of the Periodic Table. Compared with Group I metals, they:

\begin{itemize}
\item have high densities;

\item have high melting points;

\item form \textbf{coloured} 有色 \textbf{compounds} 化合物;

\item often act as catalysts, both as elements and in compounds.

\end{itemize}

They can also have \textbf{variable} 可变 \textbf{oxidation numbers} 氧化数. For example, \textbf{iron} 铁 forms both iron(II) and iron(III) compounds.

\section*{Group VIII — the noble gases}

The Group VIII \textbf{noble gases} 稀有气体 are \textbf{unreactive} 不活泼, \textbf{monatomic} 单原子 gases (they exist as single atoms, not as molecules).

They are unreactive because they already have a full outer shell of electrons. This makes them stable, so they do not need to gain, lose or share electrons.


\end{document}
